\documentclass[11pt]{article}
\usepackage[utf8]{inputenc}
\usepackage[T1]{fontenc}
\usepackage[a4paper, margin=1in]{geometry}
\usepackage{setspace}
\usepackage{microtype}
\usepackage{lmodern}

\pagestyle{empty}
\setlength{\parindent}{0pt}

\begin{document}

\begin{center}
{\Large\bfseries Trade Exposure and the Persistence of Petty Production\\[2pt]
in Post-Liberalization India\par}

\vspace{14pt}

Kabeer Bora\\[2pt]
\textit{University of Missouri--Kansas City}

\vspace{18pt}

\textbf{ABSTRACT}
\end{center}

\vspace{4pt}

Does trade integration transform how firms are socially organized? In India, seven out of ten establishments still hire no workers. This has persisted ever since data on unincorporated establishments were first collected and released. Standard trade theory, along with the misallocation literature, would say that trade should reduce misallocation, thereby removing the most unproductive firms, which are these family-run firms with no hired labor. This is an empirical question, which we put to the test in this paper by asking whether the share of these firms has decreased due to increasing trade liberalization in India. Using the universe of firms in the Economic Census from 1998 to 2015 and a measure of each district's exposure to world demand, we find that the effect of trade liberalization on this composition is precisely null. Other work on formalization has found that the number of firms hiring formal workers increases due to trade liberalization; however, the evidence suggests that the social organization of work is largely untouched by trade, both in manufacturing establishments and across non-farm employment.

\end{document}
